\begin{proposition}[Proved: the natural-order pullback bases do not lift to words]
\label{G:natural-order-fiber-obstruction}
For the independently ordered punctured $\mathbb F_2$-vector-space
category, let
\[
 N(B)=\mathbb Z^{\mathcal T(B)}
\]
be the natural-order function diagram from
Proposition~\ref{A:ordered-vector-integral-torsion}, with its
point-indicator bases. There is no free-monoid diagram with these
specified free bases whose abelianization maps are the pullback
maps of $N$. The obstruction already occurs on one finite
commuting diagram with objects of dimensions one, two, and three.
This assertion does not exclude a different choice of bases or a
stabilized positive realization of a different diagram.
\end{proposition}
\begin{proof}
Let $V$ be the punctured one-dimensional space, let $P$ be the
punctured two-dimensional space ordered $x_1<x_2<x_3$, and let
$Q$ be the punctured three-dimensional space with its binary labels
in the reference order $1<2<\cdots<7$. Define
\[
 a(P)=(1,2,3),\qquad b(P)=(1,4,5),
 \qquad v:V\to P,\quad v(V)=x_1.
\]
These are increasing linear embeddings and $av=bv$. The sets
of auxiliary natural orders have respective sizes $1,6,168$.
In the point-indicator bases the pullback of one order of $P$ is
the sum of all $28$ orders of $Q$ restricting to it.

Any word with this abelianization must list exactly those $28$
distinct generators, once each. Similarly the image of the sole
generator at $V$ in the free monoid at $P$ must list all six
generators once. Functoriality along $av=bv$ would therefore give
one word listing the $168$ generators at $Q$, in which every fiber
of either restriction $\mathcal T(Q)\to\mathcal T(P)$ is a
contiguous block. Thus the two six-part restriction partitions
would both be interval partitions of one linear order.

Two interval partitions with six nonempty blocks each have at
most $6+6-1=11$ nonempty pairwise intersections: their combined
ten interior cuts divide the line into at most eleven intervals.
Here all $36$ intersections are nonempty. In the six order types
\[
 012,\ 021,\ 102,\ 120,\ 201,\ 210
\]
on the locally indexed points of each plane, exact enumeration of
the $168$ ordered bases gives the joint fiber cardinalities
\[
 \begin{pmatrix}
 10&10&2&2&2&2\\
 10&10&2&2&2&2\\
 2&2&8&4&8&4\\
 2&2&4&8&4&8\\
 2&2&8&4&8&4\\
 2&2&4&8&4&8
 \end{pmatrix}.
\]
Every entry is positive and every row and column sums to $28$.
The attached exhaustive script constructs each natural order by
an ordered basis and verifies this matrix with integer arithmetic.
The interval-partition bound is contradicted, proving the claim.
\end{proof}
